% Abstract body only — \input inside the abstract environment in main.tex.
Work-related accidents kill nearly three million people each year and cost an estimated 4\% of global GDP, yet the IP camera networks already installed across industrial sites still depend on human observers whose sustained attention collapses within 20--35 minutes. Automated PPE detection could close that gap, but published systems are trained and evaluated on construction-centric benchmarks and, to our knowledge, none has been evaluated on imagery from a real, operating chemical plant. This paper presents a PPE compliance pipeline developed for and evaluated at OCP Group phosphate-processing facilities (specifically at Safi's Maroc Chimie chemical plant). A YOLOv8n detector (3.0M parameters) is fine-tuned at $640\times640$ resolution on a composite industrial dataset (the public Roboflow PPE-Detection corpus, SH17 dataset (capped), the Roboflow gas masks corpus, and weighted site-specific OCP imagery) that covers the gas masks, IP camera viewing angles, and site uniforms construction benchmarks omit. The central experiment compares this industrial-adapted model against a detector trained only on the standard construction benchmark: on held-out chemical-plant imagery the adapted model reaches 65.5\% mAP@50 against 18.5\% for the generic baseline, with a wider gap still on an industrial proxy test set (49.6\% against 0.4\%). The gap is evidence that generic benchmark training does not transfer to operating-plant conditions, and that the contribution is domain-specific rather than architectural. The pipeline requires no new camera hardware and retrofits onto the IP camera infrastructure a site already operates.
